% DELTA23 characters workstream, 2026-11-21 Europe/Berlin.
% D23.C4 is an arithmetic derivation. D23.C5 was expanded after a root
% review to separate ordinary Legendre paths from the projective orbifold.
% Root independently read and passed C4--C7 on 2026-11-21, 00:29 CET;
% see the root check register for its exact event time.

\subsection{The arithmetic size of the normalized adelic frame target}

Put $O_3=\mathbb Z[\zeta_3]$, $O_4=\mathbb Z[i]$,
$K_j=O_j\otimes\mathbb Q$, and write a bar for the nontrivial
automorphism of $K_j$.  Set
\[
 N_j=\{u\in(O_j\otimes\widehat{\mathbb Z})^\times:u\bar u=1\},
 \qquad N_j(4)=\{u\in N_j:u\equiv1\pmod4\}.
\]
Here the congruence modulo four concerns the component over two.
The integral matrix embeddings are the ones in DELTA\_ON\_GT
AD.2 (integrated A.72), so $N_3(4)$ and $N_4(4)$ are precisely
its two normalized targets.

\paragraph{D23.C4 (arithmetic derivation; first proof reading passed).}
For $E_3:y^2=x^3+1$ and $E_4:y^2=x^3-x$, choose CM-compatible
bases of their full Tate modules.  Let $\rho_j$ be their adelic
representations and set $G^{\mathrm{cyc}}=G_{\mathbb Q^{\rm ab}}$.
Then
\[
 \rho_4(G^{\mathrm{cyc}})=N_4(4).
\]
For $r_0=(3/4)^{1/4}$ and
$\kappa(\sigma)=\sigma(r_0)/r_0\in\{\pm1\}$ on
$G^{\mathrm{cyc}}$, the normalized representation
\[
 \rho_3^\circ(\sigma)=\kappa(\sigma)\rho_3(\sigma)
\]
has image $N_3(4)$ and has the same kernel as $\rho_3$ on
$G^{\mathrm{cyc}}$.  Both full division fields are
$\mathbb Q(E_j[\mathrm{tors}])=K_j^{\rm ab}$.

\emph{Proof.}
We use the main theorem of complex multiplication in the torsion-point
form: Milne, \emph{Modular Functions and Modular Forms}, version 1.10,
Theorem 12.27, printed p.124
(\url{https://www.jmilne.org/math/CourseNotes/MF110.pdf}).
For a class-number-one imaginary quadratic ring $O$, it states that
the image of the full torsion representation after quotient by the
geometric automorphism group $O^\times$ is all of
$\widehat O^\times/O^\times$, with the usual ray-class interpretation.
It also implies that adjoining the Weber values of all torsion points
generates $K^{\rm ab}$.  All torsion coordinates are abelian over $K$
because the CM action is defined there and its Tate-module centralizer
is the commutative group $\widehat O^\times$.  Hence the full division
field over $K$ is $K^{\rm ab}$.

Let $\pi_4=1+i$ and $\pi_3=\sqrt{-3}$.  The entire subgroup
$E_4[\pi_4^3]$ is $K_4$-rational: besides the four two-torsion
points, the four points with abscissae $i,-i$ have ordinates
$\pm(1-i),\pm(1+i)$, respectively, and are the remaining points
of this kernel.  This last kernel identification can be checked
using the automorphism $[i](x,y)=(-x,iy)$ and the chord-and-tangent
addition law: applying $1+i$ to any of those four points gives
one of the nonzero two-torsion points in $E_4[\pi_4^2]$ which
maps to $(0,0)$ under $1+i$.
Thus $\rho_4(G_{K_4})$ is contained in
\[
 C_4=\{u\in\widehat O_4^\times:u\equiv1\pmod{\pi_4^3}\}.
\]
The reduction $\mu_4\to(O_4/\pi_4^3)^\times$ is an isomorphism:
the target has four elements, and the four roots of unity are distinct
modulo $\pi_4^3$.  Consequently $C_4\to\widehat O_4^\times/\mu_4$
is an isomorphism.  The main theorem therefore gives equality
$\rho_4(G_{K_4})=C_4$.

For $E_3$, all two-torsion points have abscissae
$-1,-\zeta_3,-\zeta_3^2$, and the kernel of $1-\zeta_3$ is
$\{O,(0,1),(0,-1)\}$.  The coprime kernels show that the whole
$E_3[2\pi_3]$ is $K_3$-rational.  The reduction
\[
 \mu_6\longrightarrow (O_3/(2\pi_3))^\times
       \simeq\mathbb F_4^\times\times\mathbb F_3^\times
\]
is an isomorphism: the cubic roots give the first factor and the
sign gives the second.  The same argument proves
\[
 \rho_3(G_{K_3})=C_3
   =\{u\in\widehat O_3^\times:u\equiv1\pmod{2\pi_3}\}.
\]
The CM fields themselves occur in the torsion fields (use the
three roots of the two-division polynomial for $E_3$, and the
Weil pairing on four-torsion for $E_4$), proving the asserted
identities over $\mathbb Q$ as well.

Complex conjugation acts on each Cartan by $u\mapsto\bar u$.
Since $C_j$ is abelian and the quotient of the full Galois image
by $C_j$ has order two, its closed derived subgroup is
$\{u/\bar u:u\in C_j\}$.  This image is compact, so no further
closure changes it.  Kronecker--Weber identifies
$G^{\mathrm{cyc}}$ with the closed derived subgroup of
$G_{\mathbb Q}$; taking the image of a closed derived subgroup
under a compact surjection gives the closed derived subgroup of
the target.  Therefore
\[
 \rho_j(G^{\mathrm{cyc}})=\{u/\bar u:u\in C_j\}.
\]

At every unramified finite place, the map $u\mapsto u/\bar u$
from local units is onto the norm-one local units.  For a split
place this is the explicit map $(a,b)\mapsto(a/b,b/a)$.
For an unramified quadratic local field extension, Hilbert 90
provides a preimage, and multiplying by a base-field uniformizer
power makes it a unit because the valuations have the same value
group.  At the inert prime two of $K_3$, a norm-one unit equal to
one modulo two has a Hilbert-90 unit preimage whose residue is
fixed by the residue-field involution.  That residue is in
$\mathbb F_2^\times=\{1\}$, so the preimage already lies in
$1+2O_{3,2}$.  Conversely such a preimage has a ratio equal to
one modulo two.

At the ramified prime three of $K_3$, a norm-one unit $v$ has
residue $\pm1$ modulo $\pi_3$.  If $v\equiv1$, a Hilbert-90
preimage has even $\pi_3$-valuation: a preimage of odd valuation
has ratio with residue $\pi_3/\bar\pi_3=-1$.  Multiplying by
a rational power of three makes the preimage a unit; multiplying
by $1$ or $-1$ makes its residue one.  Conversely a unit preimage
has ratio congruent to one.  Hence
\[
 \rho_3(G^{\mathrm{cyc}})
 =\{u\in N_3:u_2\equiv1\pmod2,
                     \ u_3\equiv1\pmod{\pi_3}\}=:M_3.
\]

At the ramified prime two of $K_4$, a unit
$u=a+bi\equiv1\pmod{\pi_4^3}$ has even $b$, and therefore
$u/\bar u\equiv1\pmod4$.  Conversely let $v\in N_{4,2}$
be one modulo four.  Put $b_0=(1+v)/2$.  It is a unit, is one
modulo two, and $b_0/\bar b_0=v$.  Among the four residue classes
of units modulo $\pi_4^3$, the two equal to one modulo two are
$1$ and $-1$.  Thus one of $b_0,-b_0$ belongs to
$1+\pi_4^3O_{4,2}$ and has the same ratio.  This proves
$\rho_4(G^{\mathrm{cyc}})=N_4(4)$.

The norm-one subgroup of $(O_3/4)^\times$ which is one modulo
two is $\{1,-1\}$.  For example write a representative as
$1+2a$; its norm is one modulo four exactly when the trace of
$a\in\mathbb F_4$ is zero, i.e. when $a=0$ or $1$.
By DELTA\_ON\_GT A.13, or its explicit four-division formulas,
the sign of $\rho_3(\sigma)$ modulo four is exactly
$\kappa(\sigma)$.  Thus $\rho_3^\circ$ takes values in $N_3(4)$.

The map $M_3\to N_3(4)$ which multiplies by this modulo-four
sign is an isomorphism.  Its inverse sends $v\in N_3(4)$ to
$\varepsilon(v_3)v$, where $\varepsilon(v_3)\in\{\pm1\}$ is
the residue of its three-component modulo $\pi_3$.  Indeed this
product has three-component equal to one modulo $\pi_3$, and
two-component equal to that sign modulo four.  Both compositions
are identity.  This proves surjectivity and the assertion about
the kernel of the normalized representation.\hfill$\square$

\paragraph{D23.C5 (proved; uniform signed integral path comparison).}
In the positive normalization of LM7\_PUSH E.65 and
DELTA\_ON\_GT AD.2, for every
$\sigma\in G_{\mathbb Q^{\rm ab}}$ one has, in CM-compatible
transported bases,
\[
 \mathcal W_\sigma=\rho_4(\sigma),\qquad
 \mathcal T_\sigma=\kappa(\sigma)\rho_3(\sigma).
\]
These are full adelic identities, not merely identities modulo twelve.

\emph{Proof, with the path marking and signed normalization made explicit.}
Use the Legendre family
$\mathscr E_t:Y^2=X(X-1)(X-t)$ on
$Y=\mathbb P^1\setminus\{0,1,\infty\}$, with the fixed positive
peripheral generators and the rational tangent $\overrightarrow{01}$.
The embedded Galois action on $F$ is this normalized tangential
action.  Its geometric monodromy in a symplectic
basis is $x\mapsto U^2$, $y\mapsto L^2$: a positive loop about
either collision of two branch points is the square of the half
twist, whose lift to the double cover is the corresponding Dehn
twist, and the two vanishing cycles have intersection one.  This
gives the signed free-subgroup representation used in AD.2.

At the tangential cusp over $t=0$, the Tate degeneration has a
primitive rank-one vanishing submodule and a rank-one quotient.
The special cubic is $Y^2=X^2(X-1)$; its two branches at the
node have slopes $\pm i$.  Thus its only splitting character is
the quadratic character of $\mathbb Q(i)$, and this is trivial on
$G_{\mathbb Q^{\rm ab}}$.  The vanishing line has cyclotomic
action and the quotient has trivial action.  Both therefore have
action one on that subgroup, including integrally at two.
The actual arithmetic matrix $V^{\rm geom}_\sigma$ at the cusp
is consequently upper unipotent with lower-right entry one.
Functoriality gives
$q\Phi_\sigma=\operatorname{Ad}(V^{\rm geom}_\sigma)q$ on $F$.
This already identifies it with $\mathcal V_\sigma$: an upper
unipotent matrix commuting with $U^2,L^2$ is the identity over
every $\mathbb Z_\ell$, including $\ell=2$, because those rings
are domains.  The normalized intertwiner of AD.1 has the same
covariance on $F$, so $V^{\rm geom}_\sigma=\mathcal V_\sigma$.

Here is the independent orbifold identification of the paths.
The quotient stack $[Y/S_3]$ has geometric fundamental group
$B=C_2\amalg C_3$ and pure subgroup $F$.  The transposition
interchanging $0$ and $1$ fixes $p=1/2$; its isotropy is the
class of $s$.  Choose an ordinary geometric path $\xi$ in $Y$
from the tangential base point to $p$ so that the transported
orbifold isotropy is the named $s$.  Such a choice exists: all
order-two factor subgroups are conjugate in $B$, and conjugators
preserving the prescribed transposition in $S_3$ lie in
$F\langle s\rangle$, so the conjugation may be made pure.
For the order-three point one uses one of
$p=(1\pm\sqrt{-3})/2$ and its positive local isotropy, and makes
the same choice with the named $r$.  These two points account
for the two possible pure conjugacy classes; select the one
corresponding to $r$.  Both points and their coordinate-permutation
isotropies are fixed by $G_{\mathbb Q^{\rm ab}}$.

Paths are composed as maps of fibre functors: $\xi$ goes from
the base point to $p$, so
\[
 c_\sigma=(\sigma\xi)^{-1}\xi\in F
\]
is a loop at the base point.  If $s_p$ denotes the isotropy at
$p$ and $s=\xi^{-1}s_p\xi$, then
\[
 \Phi_\sigma(s)=(\sigma\xi)^{-1}s_p(\sigma\xi)
                  =c_\sigma s c_\sigma^{-1}.
\]
The same formula holds with $r$.  The actual arithmetic action
on this orbifold is the source's normalized extension: it agrees
on $F$ with the given tangential action, fixes $\tau$ when the
multiplier is one, and its extension is unique since the
centralizer of $F$ in $B$ is trivial.  Uniqueness of the pure
torsion frames therefore gives $c_\sigma=\beta_\sigma$ or
$c_\sigma=A_\sigma$, respectively.

For the remaining linear calculation use the Tate sheaf only on
$Y$, where the Legendre family is an actual family; no descent of
that sheaf through the projective modular stack is asserted.
If $\rho_p(\sigma)$ is its arithmetic matrix at the ordinary
point $p$, functoriality on the path gives
\[
 q(\sigma\xi)=\rho_p(\sigma)q(\xi)\mathcal V_\sigma^{-1}.
\]
Define $M_p=q(\xi)^{-1}\rho_p(\sigma)q(\xi)$ in the
transported integral basis.  Taking the inverse and multiplying
by $q(\xi)$ shows, with the order fixed,
\[
 q(c_\sigma)=\mathcal V_\sigma M_p^{-1}.
\]
Consequently
\[
 q(\beta_\sigma)M_p(\sigma)=\mathcal V_\sigma
 \quad\hbox{or}\quad
 q(A_\sigma)M_p(\sigma)=\mathcal V_\sigma.
\]
The matrices $M_p$ are exactly $\mathcal W_\sigma$ and
$\mathcal T_\sigma$.  The path cocycle is an actual element of
$F$, so its signed value is the already specified representation
$x\mapsto U^2,y\mapsto L^2$.  There is no independent scalar
lift to choose.  The calculation holds on every finite torsion
sheaf with compatible transition maps, hence on the full integral
Tate module at every prime.

The harmonic fixed point for the interchange is $t=1/2$.
Putting $X=X'+1/2$ in its equation gives
$Y^2=X'^3-X'/4$, the quadratic twist of $E_4$ by $1/2$.
Its twisting square root belongs to $\mathbb Q^{\rm ab}$, so
its representation on $G_{\mathbb Q^{\rm ab}}$ is $\rho_4$.
For the equianharmonic point $t=(1+\sqrt{-3})/2$, since
$t^2-t+1=0$, the substitution
$X=X'+(1+t)/3$ gives
\[
 Y^2=X'^3+\frac{\sqrt{-3}}9
      =X'^3+\left(\frac1{\sqrt{-3}}\right)^3.
\]
This is the quadratic twist of $E_3$ by $d=1/\sqrt{-3}$.
If a square root of $d$ is chosen, its quotient by $r_0$ lies
in $\mathbb Q^{\rm ab}$ (its square is $-2i/3$, and one may
use $\sqrt{2/3}\,\zeta_8^{-1}$).  Thus its twisting character
on $G_{\mathbb Q^{\rm ab}}$ is $\kappa$, as claimed.
The conjugate equianharmonic point changes the twist by $-1$,
whose square root lies in $\mathbb Q^{\rm ab}$, and yields the
same conclusion.  Choose the transported CM identifications
so that their centralizers are those of $S$ and $R_0$, in the
signed convention of AD.2.  This gives the asserted exact
integral matrix identities in those bases.\hfill$\square$

\paragraph{D23.C6 (proved from D23.C4--C5 and the stipulated inputs).}
The signed adelic frame pair has exact images
\[
 (\mathcal T,\mathcal W)(H)=N_3(4)\times N_4(4),
\]
\[
 (\mathcal T,\mathcal W)(G_{\mathbb Q^{\rm ab}})
 =\{(t,w):\epsilon_3(t\bmod3)=\epsilon_2(w\bmod3)\}.
\]
Every continuous relative abelian character of $G$ whose
restriction to $H$ factors through the full adelic frame pair
factors through $\Delta$.  In particular its quadratic possibilities
are exactly $0$ and $\Delta$.

\emph{Proof.}
D23.C4 and D23.C5 show that the arithmetic pair projects onto
both normalized tori, and that their kernels on
$G_{\mathbb Q^{\rm ab}}$ cut out $K_3^{\rm ab}$ and
$K_4^{\rm ab}$.  DELTA\_ON\_GT A.12 gives their exact
intersection
\[
 K_3^{\rm ab}\cap K_4^{\rm ab}
       =\mathbb Q^{\rm ab}(3^{1/4}),
 \qquad [\mathbb Q^{\rm ab}(3^{1/4}):\mathbb Q^{\rm ab}]=2.
\]
The compact Galois image is therefore the fibre product over a
quotient of order two, and has index two in the torus product.
E.65's already checked comparison modulo twelve puts it in the
equal-parity subgroup displayed above.  That subgroup itself
has index two, so they are equal.

AD.2 puts the full gentle frame image inside the normalized torus
product.  It contains the arithmetic index-two subgroup and the
E.58 element, whose parity pair is $(a_3,a_2)=(1,0)$.
A subgroup containing an index-two subgroup and one element
outside it is the entire group.  This proves the first equality.
The quotient of the full frame image by the arithmetic image is
thus $C_2$, with quotient character $a_2+a_3$.  Any continuous
abelian character satisfying the factorization and arithmetic
vanishing hypotheses descends to that quotient.  Finally
$G=HG_{\mathbb Q}$ makes its full extension unique, and E.65
identifies the nonzero quadratic one with $\Delta$.
\hfill$\square$

\paragraph{D23.C7 (proved; precise failure of a place-two frame factorization).}
There is an actual $\sigma\in G_{\mathbb Q^{\rm ab}}$ with
\[
 a_2(\sigma)=a_3(\sigma)=1,
 \qquad
 q_2(A_\sigma)=q_2(\beta_\sigma)
     =\mathcal V_{\sigma,2}
     =\mathcal T_{\sigma,2}
     =\mathcal W_{\sigma,2}=I.
\]
Here $q_2$ means the full $2$-adic representation, not reduction
modulo two.  Thus neither parity is a function solely of the pair
of raw signed $2$-adic frame matrices, or of the corrected pair,
even if the $2$-adic intertwiner is also supplied.  This is a
statement about those precise congruence data; it makes no claim
about every possible meaning of a ``local'' geometric invariant.

There is also an actual $h\in H$
with $\Delta(h)=1$, $a_3(h)=0$, and $a_2(h)=1$, whose entire
corrected adelic pair has
\[
 \mathcal T_h=I,\qquad
 \mathcal W_{h,3}=S,\qquad
 \mathcal W_{h,\ell}=I\quad(\ell\ne3).
\]
In particular its corrected $2$-adic pair is the identity pair.

\emph{Proof.}
Fix an embedding into a decomposition group at three and write
$I_3$ for inertia.  Its wild subgroup is pro-three and its tame
quotient is $\prod_{\ell\ne3}\mathbb Z_\ell(1)$.  A pro-two
Sylow subgroup $P$ therefore maps isomorphically to the
$\mathbb Z_2(1)$ factor: its intersection with the wild subgroup
is trivial, and the image of a Sylow subgroup under a compact
surjection is a Sylow subgroup.  These assertions can also be
obtained by taking compatible Sylow subgroups in the finite
inertia quotients.  Choose a topological generator $\tau$ of $P$.
The extension generated by $r_0$ is tamely ramified of degree
four after adjoining $\mu_4$, which is unramified at three;
indeed $(2r_0)^4=12$ is Eisenstein at three.  It follows that
$\tau(r_0)=i^{\pm1}r_0$.  Put $\sigma=\tau^2$.  Then
$\sigma(r_0)=-r_0$.

For every $\ell\ne3$, the $\ell$-adic cyclotomic character is
unramified at three.  The image of the pro-two group $P$ under
$\chi_3$ is contained in $\{1,-1\}\subset\mathbb Z_3^\times$,
so its square is one.  Consequently $\chi(\sigma)=1$ and
$\sigma\in G_{\mathbb Q^{\rm ab}}$.  R4's exact finite arithmetic
comparison identifies both parities with the sign on $r_0$;
therefore $a_2(\sigma)=a_3(\sigma)=1$.

The two Legendre fibres used in the path proof of D23.C5 have
good reduction at three.  The harmonic fibre at $t=1/2$ has
discriminant $16t^2(1-t)^2=1$.  For either equianharmonic fibre,
$t(1-t)=1$, so the discriminant is $16$, a unit at three.
The latter fibre is defined over $K_3$, which $\sigma$ fixes;
$\sigma\in I_3\cap G_{K_3}$ is in its local inertia group.
The good-reduction implication of Serre--Tate,
\emph{Good reduction of abelian varieties}, Theorem 1,
\emph{Annals of Mathematics} 88 (1968), p.493, makes both full
$2$-adic Tate actions of $\sigma$ the identity.  The integral
path identity in D23.C5 therefore gives
$\mathcal T_{\sigma,2}=\mathcal W_{\sigma,2}=I$.

The actual tangential action on the free pro-two group is
unramified at three.  One explicit source for exactly this
automorphism statement is Hain--Matsumoto,
\emph{Weighted completion of Galois groups and Galois actions on
the fundamental group of $\mathbb P^1-\{0,1,\infty\}$},
\emph{Compositio Mathematica} 139 (2003), Section 8.5,
the paragraph before Theorem 8.7 (take $N=1$ and $\ell=2$).
Thus $\Phi_\sigma$ acts identically on $F^{(2)}$.
The image $q_2(F)$ lies in the principal congruence subgroup
modulo two of $\mathrm{SL}_2(\mathbb Z_2)$, which is pro-two.
Hence $q_2\Phi_\sigma=q_2$.  The covariance says that the
upper unipotent $\mathcal V_{\sigma,2}$ commutes with $U^2,L^2$,
forcing it to be $I$.  The corrected-frame definitions then give
$q_2(A_\sigma)=q_2(\beta_\sigma)=I$ as well.

The identity element has all those same matrices and parity zero,
which proves the asserted failure of factorization.  For the final
assertion, the displayed adelic pair belongs to
$N_3(4)\times N_4(4)$: its component over two is the identity,
and $S$ has norm one at three.  D23.C6 supplies an actual $h$
with that value.  Its parities and discrepancy have the stated
values by their definitions.\hfill$\square$

\paragraph{D23.C8 (proved; an explicit smaller gate for B.40).}
Let
\[
 N=\ker\bigl((\mathcal T,\mathcal W):H\longrightarrow
                   N_3(4)\times N_4(4)\bigr),
 \qquad
 Q=G/\overline{[G,G]G^2G_{\mathbb Q}}.
\]
Then, as a set product,
\[
 \ker\Delta=N G_{\mathbb Q}.
\]
If $\bar\Delta:Q\to\mathbb F_2$ is the induced nonzero
character, its kernel is exactly the image of $N$.  Consequently
the assertion that all relative continuous quadratic characters
are $0,\Delta$ is equivalent to the concrete vanishing condition
\[
 N\subseteq\overline{[G,G]G^2G_{\mathbb Q}}.
\]
This computes the entire remaining kernel in terms of elements
whose two full adelic frames are identity.  It does not assert
that this subgroup has trivial image in $Q$.

For the two particular characters $d$ and $\Delta\circ D$,
the known actual kernel witness of D23.C3 gives the further
equivalences
\[
 \begin{split}
 d\ne0&\quad\Longleftrightarrow\quad
      \exists n\in N:\ a_2(Dn)=1,\\
 \Delta\circ D\ne0&\quad\Longleftrightarrow\quad
      \exists n\in N:\ \Delta(Dn)=1.
 \end{split}
\]
In particular both characters vanish if and only if
$(a_2,a_3)(D(N))=\{(0,0)\}$.

\emph{Proof.}
Write $g=h\sigma$ with $h\in H$ and $\sigma\in G_{\mathbb Q}$.
If $g\in\ker\Delta$, then $h\in\ker\delta$, so its two
adelic frame values have equal parity.  D23.C6 gives an actual
$\tau\in G_{\mathbb Q^{\rm ab}}$ with exactly the same entire
frame pair as $h$.  The two frame maps are homomorphisms on
$H$, hence $n=h\tau^{-1}\in N$ and $g=n(\tau\sigma)$.
The converse follows from $\Delta|_N=\Delta|_{G_{\mathbb Q}}=0$.
No normality assertion about $N$ in all of $G$ is needed.

The closed subgroup in the definition of $Q$ is the inverse
image of the image of $G_{\mathbb Q}$ in the maximal elementary
abelian pro-two quotient, so the quotient is well defined.
The preceding set equality shows
$\ker\bar\Delta=\operatorname{im}(N\to Q)$.  Continuous
quadratic characters separate points of an elementary abelian
pro-two group.  Thus they are exactly $0,\bar\Delta$ if and
only if this kernel is trivial.  Pullback gives the asserted
equivalence on $G$.  In the integrated relative-abelianization
statement P.12, $Q=K_G/2K_G$.

On $N$, one has $a_2=a_3=0$, so
$d(n)=a_2(Dn)$.  Any relative character vanishing on $N$
vanishes on $\ker\Delta=N G_{\mathbb Q}$ and is therefore
$0$ or $\Delta$.  The actual $k$ of D23.C3 has
$d(k)=(\Delta\circ D)(k)=0$ but $\Delta(k)=1$;
it excludes the second possibility for both characters.
This proves the two equivalences.  Their right sides both
fail exactly when $a_2(Dn)=0$ and
$a_2(Dn)+a_3(Dn)=0$ for every $n\in N$, equivalently when
both descended parities vanish there.\hfill$\square$
